% GENERATED FILE. Edit canonical lessons, then run npm run generate:books.
% canonical-source-hash: fnv1a64:da16cddb566cdfff
% canonical-lessons: TA-C454-tikil, TA-C454-piti, TA-C454-atircci, TA-C454-titukkitu, TA-C454-natunku

\chapter{Fear and Shock}
\label{ch:ta-fear-and-shock}
\hlchaptermodality{454}

\begin{chapteropening}
\textbf{By the end of this chapter:} \emph{I can talk about terror, panic and being startled.}

Complete the last lesson of chapter 454: I can talk about terror, panic and being startled.
\end{chapteropening}

\section[\texorpdfstring{\ta{திகில்} (tikil)}{திகில் (tikil)}]{\ta{திகில்} (tikil) — terror; horror}

\label{lesson:TA-C454-tikil}



\begin{quote}
Before the new one: say the Tamil for to seethe, then the Tamil for to mutter.
\end{quote}

\subsection*{You'll want to know: \ta{திகில்}}
\textbf{\ta{திகில்}} — \emph{tikil} — \textquotedblleft{}terror; horror\textquotedblright{}.

\begin{quote}
\ta{அந்தக்} \ta{கதையைக்} \ta{கேட்டு} \ta{எனக்குத்} \ta{திகிலாக} \ta{இருந்தது}.
\end{quote}

\begin{quote}
\emph{antak kataiyaik kēṭṭu eṉakkut tikilāka iruntatu.}
\end{quote}

\begin{quote}
Hearing that story, I felt terrified.
\end{quote}

With \ta{இரு}: \ta{திகிலாக} \ta{இருந்தது}, it was terrifying. A horror film is a \ta{திகில்} \ta{படம்}.

\begin{cousinweb}
\textbf{\ta{திகில்}} (\emph{tikil}) is native Dravidian (DEDR 3202); Kota in the same entry has a verb for being thunderstruck.
\end{cousinweb}

\begin{grammarlens}[title={What you've built}]
One more word for describing what you lived and felt.
\end{grammarlens}

\subsection*{Your turn}
\begin{itemize}
\raggedright
  \item \emph{Say it:} \emph{tikil}
  \item \emph{Say it:} \emph{tikil}, once more
  \item \emph{Say it:} \emph{antak kataiyaik kēṭṭu eṉakkut tikilāka iruntatu.} — Hearing that story, I felt terrified.
  \item \emph{Recall:} say the Tamil for to seethe, then the Tamil for to mutter, then say \emph{tikil} again
\end{itemize}

\begin{tcolorbox}[breakable,colback=teal!4,colframe=teal!35!black,title={Before you move on}]
What does \ta{திகில்} mean? (\textquotedblleft{}Terror; horror\textquotedblright{}.) Say it once more.
\end{tcolorbox}
\section[\texorpdfstring{\ta{பீதி} (pīti)}{பீதி (pīti)}]{\ta{பீதி} (pīti) — panic; dread}

\label{lesson:TA-C454-piti}



\begin{quote}
Before the new one: say the Tamil for to mutter, then the Tamil for terror.
\end{quote}

\subsection*{You'll want to know: \ta{பீதி}}
\textbf{\ta{பீதி}} — \emph{pīti} — \textquotedblleft{}panic; dread\textquotedblright{}.

\begin{quote}
\ta{நிலநடுக்கத்தின்போது} \ta{மக்களிடம்} \ta{பீதி} \ta{பரவியது}.
\end{quote}

\begin{quote}
\emph{nilanaṭukkattiṉpōtu makkaḷiṭam pīti paraviyatu.}
\end{quote}

\begin{quote}
During the earthquake panic spread among the people.
\end{quote}

\ta{பீதி} \ta{அடை}, to panic (literally to reach panic).

\begin{cousinweb}
\textbf{\ta{பீதி}} (\emph{pīti}) is borrowed from Sanskrit \emph{bhīti}, \textquotedblleft{}fear\textquotedblright{} (CDIAL 9512); Tamil writes the \emph{bh} as \ta{ப}.
\end{cousinweb}

\begin{grammarlens}[title={What you've built}]
One more word for describing what you lived and felt.
\end{grammarlens}

\subsection*{Your turn}
\begin{itemize}
\raggedright
  \item \emph{Say it:} \emph{pīti}
  \item \emph{Say it:} \emph{pīti}, once more
  \item \emph{Say it:} \emph{nilanaṭukkattiṉpōtu makkaḷiṭam pīti paraviyatu.} — During the earthquake panic spread among the people.
  \item \emph{Recall:} say the Tamil for to mutter, then the Tamil for terror, then say \emph{pīti} again
\end{itemize}

\begin{tcolorbox}[breakable,colback=teal!4,colframe=teal!35!black,title={Before you move on}]
What does \ta{பீதி} mean? (\textquotedblleft{}Panic; dread\textquotedblright{}.) Say it once more.
\end{tcolorbox}
\section[\texorpdfstring{\ta{அதிர்ச்சி} (atircci)}{அதிர்ச்சி (atircci)}]{\ta{அதிர்ச்சி} (atircci) — shock}

\label{lesson:TA-C454-atircci}



\begin{quote}
Before the new one: say the Tamil for terror, then the Tamil for panic.
\end{quote}

\subsection*{You'll want to know: \ta{அதிர்ச்சி}}
\textbf{\ta{அதிர்ச்சி}} — \emph{atircci} — \textquotedblleft{}shock\textquotedblright{}.

\begin{quote}
\ta{அந்தச்} \ta{செய்தி} \ta{எனக்குப்} \ta{பெரிய} \ta{அதிர்ச்சியாக} \ta{இருந்தது}.
\end{quote}

\begin{quote}
\emph{antac ceyti eṉakkup periya atircciyāka iruntatu.}
\end{quote}

\begin{quote}
That news was a big shock to me.
\end{quote}

From the verb \ta{அதிர்} \textquotedblleft{}to shake, to quake\textquotedblright{}: a shock is a shaking.

\begin{cousinweb}
\textbf{\ta{அதிர்ச்சி}} (\emph{atircci}) is native Dravidian, a noun from \ta{அதிர்} \emph{atir}, \textquotedblleft{}to shake, to quake\textquotedblright{} (DEDR 137), which lists it.
\end{cousinweb}

\begin{grammarlens}[title={What you've built}]
One more word for describing what you lived and felt.
\end{grammarlens}

\subsection*{Your turn}
\begin{itemize}
\raggedright
  \item \emph{Say it:} \emph{atircci}
  \item \emph{Say it:} \emph{atircci}, once more
  \item \emph{Say it:} \emph{antac ceyti eṉakkup periya atircciyāka iruntatu.} — That news was a big shock to me.
  \item \emph{Recall:} say the Tamil for terror, then the Tamil for panic, then say \emph{atircci} again
\end{itemize}

\begin{tcolorbox}[breakable,colback=teal!4,colframe=teal!35!black,title={Before you move on}]
What does \ta{அதிர்ச்சி} mean? (\textquotedblleft{}Shock\textquotedblright{}.) Say it once more.
\end{tcolorbox}
\section[\texorpdfstring{\ta{திடுக்கிடு} (tiṭukkiṭu)}{திடுக்கிடு (tiṭukkiṭu)}]{\ta{திடுக்கிடு} (tiṭukkiṭu) — to be startled}

\label{lesson:TA-C454-titukkitu}



\begin{quote}
Before the new one: say the Tamil for panic, then the Tamil for shock.
\end{quote}

\subsection*{You'll want to know: \ta{திடுக்கிடு}}
\textbf{\ta{திடுக்கிடு}} — \emph{tiṭukkiṭu} — \textquotedblleft{}to be startled\textquotedblright{}.

\begin{quote}
\ta{கதவு} \ta{திடீரென்று} \ta{திறந்தது}, \ta{நான்} \ta{திடுக்கிட்டேன்}.
\end{quote}

\begin{quote}
\emph{katavu tiṭīreṉṟu tiṟantatu, nāṉ tiṭukkiṭṭēṉ.}
\end{quote}

\begin{quote}
The door opened suddenly and I was startled.
\end{quote}

Past (I …): \ta{திடுக்கிட்டேன்}, \emph{tiṭukkiṭṭēṉ}.

Built with \ta{இடு} \textquotedblleft{}to put\textquotedblright{}, so its past runs like \ta{இட்டேன்}: \ta{திடுக்கிட்டேன்}.

\begin{cousinweb}
\textbf{\ta{திடுக்கிடு}} (\emph{tiṭukkiṭu}) is a native compound of \ta{திடுக்கு} \emph{tiṭukku}, \textquotedblleft{}a sudden fright\textquotedblright{} (DEDR 3216), and the verb \ta{இடு} \textquotedblleft{}to put\textquotedblright{}.
\end{cousinweb}

\begin{grammarlens}[title={What you've built}]
One more word for describing what you lived and felt.
\end{grammarlens}

\subsection*{Your turn}
\begin{itemize}
\raggedright
  \item \emph{Say it:} \emph{tiṭukkiṭu}
  \item \emph{Say it:} \emph{tiṭukkiṭu}, once more
  \item \emph{Say it:} \emph{katavu tiṭīreṉṟu tiṟantatu, nāṉ tiṭukkiṭṭēṉ.} — The door opened suddenly and I was startled.
  \item \emph{Recall:} say the Tamil for panic, then the Tamil for shock, then say \emph{tiṭukkiṭu} again
\end{itemize}

\begin{tcolorbox}[breakable,colback=teal!4,colframe=teal!35!black,title={Before you move on}]
What does \ta{திடுக்கிடு} mean? (\textquotedblleft{}To be startled\textquotedblright{}.) Say it once more.
\end{tcolorbox}
\section[\texorpdfstring{\ta{நடுங்கு} (naṭuṅku)}{நடுங்கு (naṭuṅku)}]{\ta{நடுங்கு} (naṭuṅku) — to tremble; to shiver}

\label{lesson:TA-C454-natunku}



\begin{quote}
Before the new one: say the Tamil for shock, then the Tamil for to be startled.
\end{quote}

\subsection*{You'll want to know: \ta{நடுங்கு}}
\textbf{\ta{நடுங்கு}} — \emph{naṭuṅku} — \textquotedblleft{}to tremble; to shiver\textquotedblright{}.

\begin{quote}
\ta{குளிரில்} \ta{குழந்தை} \ta{நடுங்கியது}.
\end{quote}

\begin{quote}
\emph{kuḷiril kuḻantai naṭuṅkiyatu.}
\end{quote}

\begin{quote}
The child shivered in the cold.
\end{quote}

Past (I …): \ta{நடுங்கினேன்}, \emph{naṭuṅkiṉēṉ}.

From cold or from fear: \ta{பயத்தில்} \ta{நடுங்கினேன்}, I trembled with fear.

\begin{cousinweb}
\textbf{\ta{நடுங்கு}} (\emph{naṭuṅku}) is native Dravidian (DEDR 3585), from cold, fear or weakness alike.
\end{cousinweb}

\begin{grammarlens}[title={What you've built}]
One more word for describing what you lived and felt.
\end{grammarlens}

\subsection*{Your turn}
\begin{itemize}
\raggedright
  \item \emph{Say it:} \emph{naṭuṅku}
  \item \emph{Say it:} \emph{naṭuṅku}, once more
  \item \emph{Say it:} \emph{kuḷiril kuḻantai naṭuṅkiyatu.} — The child shivered in the cold.
  \item \emph{Recall:} say the Tamil for shock, then the Tamil for to be startled, then say \emph{naṭuṅku} again
\end{itemize}

\begin{tcolorbox}[breakable,colback=teal!4,colframe=teal!35!black,title={Before you move on}]
What does \ta{நடுங்கு} mean? (\textquotedblleft{}To tremble; to shiver\textquotedblright{}.) Say it once more.
\end{tcolorbox}
